\documentclass[aspectratio=169, 10pt]{beamer}

% ==============================================================================
% CONFIGURACIÓN DE IDIOMA Y CODIFICACIÓN
% ==============================================================================
\usepackage[utf8]{inputenc}
\usepackage[T1]{fontenc}
\usepackage[spanish,es-tabla]{babel}

% ==============================================================================
% PAQUETES TIPOGRÁFICOS Y DE FORMATO
% ==============================================================================
\usepackage{helvet}
\renewcommand{\familydefault}{\sfdefault}
\usepackage{amsmath, amssymb, amsfonts}
\usepackage{graphicx}
\usepackage{booktabs}
\usepackage{tikz}
\usetikzlibrary{shapes.geometric, arrows.meta, positioning, calc, shadows}
\usepackage{listings}
\usepackage{tcolorbox}
\tcbuselibrary{skins,breakable}

% ==============================================================================
% DEFINICIÓN DE COLORES INSTITUCIONALES (PUCP / Q-LAB)
% ==============================================================================
\definecolor{pucpRed}{RGB}{140, 20, 40}        % Carmesí / Borgoña institucional
\definecolor{pucpBlue}{RGB}{0, 40, 85}         % Azul marino profundo
\definecolor{pucpLightBlue}{RGB}{230, 241, 252}% Azul claro para bloques y acentos
\definecolor{pucpMidBlue}{RGB}{15, 75, 140}    % Azul medio para subsecciones
\definecolor{pucpGray}{RGB}{246, 248, 250}     % Fondo claro de código y cajas
\definecolor{pucpDarkGray}{RGB}{70, 74, 78}    % Texto secundario
\definecolor{accentGreen}{RGB}{28, 125, 46}    % Verde para ejemplos / éxito
\definecolor{accentGold}{RGB}{195, 130, 0}     % Dorado para notas y tips

% ==============================================================================
% TEMA DE BEAMER Y PERSONALIZACIÓN VISUAL
% ==============================================================================
\useinnertheme{rectangles}

% Desactivar barra de navegación inferior estándar
\setbeamertemplate{navigation symbols}{}

% Colores de la estructura
\setbeamercolor{structure}{fg=pucpBlue}
\setbeamercolor{palette primary}{bg=pucpBlue, fg=white}
\setbeamercolor{palette secondary}{bg=pucpMidBlue, fg=white}
\setbeamercolor{palette tertiary}{bg=pucpRed, fg=white}
\setbeamercolor{palette quaternary}{bg=pucpBlue, fg=white}

% Encabezado (Headline) elegante con sección actual
\setbeamertemplate{headline}{%
  \leavevmode%
  \ifnum\thepage=1\else%
  \begin{beamercolorbox}[wd=\paperwidth,ht=2.4ex,dp=1.0ex,leftskip=1.2em,rightskip=1.2em]{palette primary}%
    \fontsize{7}{8}\selectfont\textbf{Q-LAB} \textbar\ \insertshorttitle \hfill \textit{\insertsectionhead}%
  \end{beamercolorbox}%
  \fi%
  \vskip0pt%
}

% Título del frame con margen interno seguro
\setbeamercolor{frametitle}{bg=pucpRed, fg=white}
\setbeamerfont{frametitle}{series=\bfseries, size=\large}
\setbeamertemplate{frametitle}{%
  \nointerlineskip%
  \begin{beamercolorbox}[wd=\paperwidth,ht=3.2ex,dp=1.2ex,leftskip=1.2em,rightskip=1.2em]{frametitle}%
    \usebeamerfont{frametitle}\insertframetitle\strut%
  \end{beamercolorbox}%
}

% Bloques estándar, de alerta y de ejemplos
\setbeamercolor{block title}{bg=pucpBlue, fg=white}
\setbeamercolor{block body}{bg=pucpLightBlue!60, fg=black}

\setbeamercolor{block title alerted}{bg=pucpRed, fg=white}
\setbeamercolor{block body alerted}{bg=pucpRed!10, fg=black}

\setbeamercolor{block title example}{bg=accentGreen, fg=white}
\setbeamercolor{block body example}{bg=accentGreen!10, fg=black}

% Listas (ítems)
\setbeamercolor{item}{fg=pucpBlue}
\setbeamercolor{subitem}{fg=pucpRed}
\setbeamertemplate{itemize item}{\color{pucpBlue}$\blacktriangleright$}
\setbeamertemplate{itemize subitem}{\color{pucpRed}$\bullet$}
\setbeamertemplate{itemize/enumerate body begin}{\small}
\setbeamertemplate{itemize/enumerate subbody begin}{\footnotesize}

% ==============================================================================
% FOOTLINE PERSONALIZADO (Estilo Q-LAB / PUCP con dimensiones precisas)
% ==============================================================================
\setbeamertemplate{footline}{%
  \leavevmode%
  \hbox{%
    \begin{beamercolorbox}[wd=0.18\paperwidth,ht=2.4ex,dp=0.9ex,center]{palette primary}%
      \textbf{Q-LAB}
    \end{beamercolorbox}%
    \begin{beamercolorbox}[wd=0.64\paperwidth,ht=2.4ex,dp=0.9ex,center]{palette secondary}%
      \scriptsize Taller Introductorio de Overleaf (\LaTeX) para la Redacción Académica
    \end{beamercolorbox}%
    \begin{beamercolorbox}[wd=0.18\paperwidth,ht=2.4ex,dp=0.9ex,center]{palette primary}%
      \textbf{PUCP}\hspace{0.6em}\insertframenumber\,/\,\inserttotalframenumber
    \end{beamercolorbox}%
  }%
  \vskip0pt%
}

% ==============================================================================
% CONFIGURACIÓN DE LISTINGS (CÓDIGO Y COMANDOS)
% ==============================================================================
\lstset{
  basicstyle=\ttfamily\scriptsize,
  backgroundcolor=\color{pucpGray},
  keywordstyle=\color{pucpRed}\bfseries,
  commentstyle=\color{pucpDarkGray}\itshape,
  stringstyle=\color{pucpMidBlue},
  numbers=none,
  breaklines=true,
  showstringspaces=false,
  tabsize=2,
  frame=single,
  rulecolor=\color{black!15}
}

% ==============================================================================
% METADATOS DE LA PRESENTACIÓN
% ==============================================================================
\title[\LaTeX\ en VS Code, Git \& IA]{Flujo de Trabajo Avanzado con \LaTeX: Integración de VS Code, Git/GitHub y Copilotos de IA}
\subtitle{Sesión 3: Del Piloto Manual al Piloto Automático}
\author[Q-LAB - PUCP]{Laboratorio de Inteligencia Artificial y Métodos Computacionales en Ciencias Sociales (Q-LAB)}
\institute[PUCP]{Pontificia Universidad Católica del Perú}
\date{26 de agosto de 2026}

% ==============================================================================
% INICIO DEL DOCUMENTO
% ==============================================================================
\begin{document}

% ------------------------------------------------------------------------------
% DIAPOSITIVA 1: PORTADA
% ------------------------------------------------------------------------------
\begin{frame}[plain]
  \vfill
  \centering
  \begin{tcolorbox}[colback=pucpRed,colframe=pucpRed,coltext=white,boxrule=0pt,arc=2.5mm,width=0.95\textwidth,halign=center,top=2.8mm,bottom=2.8mm]
    {\Large\textbf{Flujo de Trabajo Avanzado con \LaTeX}}\\[0.5ex]
    {\small\textit{Sesión 3: Integración de Overleaf, VS Code, Git/GitHub y Copilotos de IA}}
  \end{tcolorbox}
  
  \vspace{0.35cm}
  {\small\textbf{Laboratorio de Inteligencia Artificial y Métodos Computacionales en Ciencias Sociales (Q-LAB)}}\\[0.3ex]
  {\footnotesize Pontificia Universidad Católica del Perú}\\[0.5ex]
  {\scriptsize 26 de agosto de 2026}
  
  \vspace{0.35cm}
  \begin{tcolorbox}[colback=pucpBlue,colframe=pucpBlue,coltext=white,boxrule=0pt,arc=2.5mm,width=0.42\textwidth,halign=center,top=1.8mm,bottom=1.8mm]
    \textbf{\normalsize Q-LAB}\\[0.1ex]
    \tiny Innovación, Datos y Redacción Académica
  \end{tcolorbox}
  \vfill
\end{frame}

% ------------------------------------------------------------------------------
% DIAPOSITIVA 2: AGENDA / TABLA DE CONTENIDOS
% ------------------------------------------------------------------------------
\begin{frame}{Estructura y Objetivos de la Sesión 3}
  \tableofcontents
\end{frame}

% ==============================================================================
% SECCIÓN 1: INTRODUCCIÓN Y ENFOQUE PEDAGÓGICO
% ==============================================================================
\section[Piloto Automático]{De Piloto Manual a Piloto Automático}

\begin{frame}{Transición: Del Piloto Manual al Piloto Automático}
  \begin{columns}[T]
    \begin{column}{0.48\textwidth}
      \begin{block}{Sesiones 1 y 2: Piloto Manual}
        \small
        \begin{itemize}
          \item Uso exclusivo del editor web de Overleaf.
          \item Memorización manual de sintaxis y etiquetas.
          \item Flujo individual y dependiente del navegador.
          \item Restricciones de cuotas de cómputo en la nube.
        \end{itemize}
      \end{block}
    \end{column}
    
    \begin{column}{0.48\textwidth}
      \begin{block}{Sesión 3: Piloto Automático}
        \small
        \begin{itemize}
          \item Entorno de Desarrollo Integrado (\textbf{IDE: VS Code}).
          \item Control de versiones profesional (\textbf{Git y GitHub}).
          \item Asistencia y automatización con IA (\textbf{Antigravity}).
          \item Trabajo local/remoto sin límites de compilación.
        \end{itemize}
      \end{block}
    \end{column}
  \end{columns}

  \vspace{0.15cm}
  \begin{alertblock}{Objetivo Pedagógico Central}
    \small
    Aprender herramientas profesionales transferibles: no solo para redactar en \LaTeX, sino para investigación reproducible, ciencia de datos (R, Python) y colaboración multidisciplinaria.
  \end{alertblock}
\end{frame}

% ==============================================================================
% SECCIÓN 2: CONTROL DE VERSIONES CON GIT Y GITHUB
% ==============================================================================
\section[Git \& GitHub]{Control de Versiones: Git y GitHub}

\begin{frame}{¿Qué es Git y qué problema resuelve?}
  \begin{columns}[T]
    \begin{column}{0.54\textwidth}
      \begin{block}{Definición de Git}
        \small
        Sistema de control de versiones distribuido (\textit{VCS}) que registra el historial completo de cambios en archivos de texto.
      \end{block}

      \vspace{0.1cm}
      \textbf{\small ¿Cómo opera en la práctica?}
      \small
      \begin{itemize}
        \item Captura \textbf{instantáneas (\textit{commits})} en puntos clave.
        \item Permite revertir cambios o comparar versiones sin riesgo.
        \item Facilita el trabajo concurrente mediante ramas (\textit{branches}).
      \end{itemize}
    \end{column}

    \begin{column}{0.44\textwidth}
      \begin{alertblock}{El Trastorno del "final\_doc"}
        \scriptsize
        \texttt{manuscrito\_final.docx} \\
        $\to$ \texttt{manuscrito\_final\_revisado.docx} \\
        $\to$ \texttt{manuscrito\_final\_asesor.docx} \\
        $\to$ \texttt{manuscrito\_final\_ESTE\_SI.docx} \\
        $\to$ \texttt{manuscrito\_final\_v3\_definitivo.docx}
      \end{alertblock}
      \vspace{0.12cm}
      \begin{center}
        \scriptsize\textit{Git sustituye el desorden de archivos por un historial ordenado y verificable.}
      \end{center}
    \end{column}
  \end{columns}
\end{frame}

\begin{frame}{Conceptos Clave de Git Desmitificados}
  \begin{columns}[T]
    \begin{column}{0.48\textwidth}
      \begin{block}{Los 3 Estados de un Archivo en Git}
        \small
        \begin{enumerate}
          \item \textbf{Working Directory:} Carpeta local de edición.
          \item \textbf{Staging Area (\texttt{git add}):} Selección de cambios.
          \item \textbf{Repository (\texttt{git commit}):} Registro formal guardado permanentemente en el historial.
        \end{enumerate}
      \end{block}
    \end{column}

    \begin{column}{0.48\textwidth}
      \begin{block}{Colaboración con Ramas (\textit{Branches})}
        \small
        \begin{itemize}
          \item \textbf{Branch (Rama):} Línea de trabajo paralela sin alterar el documento principal (\texttt{main}).
          \item \textbf{Pull Request (PR):} Solicitud para integrar cambios.
          \item \textbf{Push / Pull:} Envío y descarga de cambios remotos.
        \end{itemize}
      \end{block}
    \end{column}
  \end{columns}

  \vspace{0.15cm}
  \begin{tcolorbox}[colback=pucpLightBlue!50,colframe=pucpBlue,boxrule=0.5pt,arc=1.5mm,top=1.5mm,bottom=1.5mm]
    \footnotesize\textbf{Analogía:} \textit{Git} es el motor local que gestiona las versiones; \textit{GitHub} es la plataforma en la nube que permite compartir y colaborar en equipo.
  \end{tcolorbox}
\end{frame}

\begin{frame}{GitHub: Colaboración, Portafolio y Educación}
  \begin{columns}[T]
    \begin{column}{0.48\textwidth}
      \begin{block}{¿Qué es GitHub?}
        \small
        \begin{itemize}
          \item Plataforma web para repositorios Git remotos.
          \item Centro de código abierto y ciencia reproducible.
          \item Permite compartir textos, datos y bibliografía.
        \end{itemize}
      \end{block}
      
      \vspace{0.1cm}
      \begin{exampleblock}{Portafolio Académico y Laboral}
        \footnotesize
        Tus repositorios sirven como carta de presentación ante asesores, revistas y empleadores.
      \end{exampleblock}
    \end{column}

    \begin{column}{0.48\textwidth}
      \begin{block}{GitHub Education (Beneficios Pro)}
        \small
        \begin{itemize}
          \item Registro con correo PUCP (\texttt{@pucp.edu.pe}).
          \item Acceso a \textbf{GitHub Pro} sin costo.
          \item Habilitación de herramientas avanzadas para investigación y docencia.
        \end{itemize}
      \end{block}

      \vspace{0.1cm}
      \begin{alertblock}{Estructura del Taller en GitHub}
        \footnotesize
        Repositorio modular organizado por carpetas (\texttt{Sesión 1}, \texttt{Sesión 2}, \dots) con código fuente y rúbricas.
      \end{alertblock}
    \end{column}
  \end{columns}
\end{frame}

% ==============================================================================
% SECCIÓN 3: VISUAL STUDIO CODE COMO ENTORNO INTEGRADO (IDE)
% ==============================================================================
\section[VS Code]{Visual Studio Code como IDE Central}

\begin{frame}{Visual Studio Code (VS Code): El Entorno Integrado}
  \begin{block}{¿Qué es un IDE (Integrated Development Environment)?}
    \small
    Un entorno que centraliza editor de texto, terminal, depurador, gestor de extensiones y visores en una sola ventana de trabajo profesional.
  \end{block}

  \vspace{0.1cm}
  \begin{columns}[T]
    \begin{column}{0.50\textwidth}
      \textbf{\small ¿Por qué centralizar en VS Code?}
      \small
      \begin{itemize}
        \item Reduce drásticamente el consumo de memoria RAM.
        \item Unifica todas tus herramientas académicas:
        \begin{itemize}\footnotesize
          \item \LaTeX\ / Overleaf
          \item Scripts en R, Python, Julia o Stata
          \item Control de versiones con Git
          \item Gestores bibliográficos (Zotero / BibTeX)
        \end{itemize}
      \end{itemize}
    \end{column}

    \begin{column}{0.48\textwidth}
      \begin{tcolorbox}[colback=pucpGray,colframe=pucpBlue,title=\textbf{Paneles Clave en VS Code},fonttitle=\bfseries\scriptsize,boxrule=0.6pt,top=1.5mm,bottom=1.5mm]
        \scriptsize
        \begin{itemize}
          \item \textbf{Explorador:} Árbol de archivos y carpetas.
          \item \textbf{Extensiones:} Instalación de plugins (Overleaf, Python).
          \item \textbf{Paleta de Comandos:} \texttt{Ctrl + Shift + P}.
          \item \textbf{Terminal Integrada:} PowerShell / Bash.
          \item \textbf{Outline (Esquema):} Índice de secciones y figuras.
        \end{itemize}
      \end{tcolorbox}
    \end{column}
  \end{columns}
\end{frame}

\begin{frame}{Terminales y Entorno: PowerShell vs. CMD vs. Bash}
  \begin{block}{La Terminal Integrada en VS Code}
    \small
    Permite interactuar directamente con el sistema operativo sin cambiar de programa.
  \end{block}

  \vspace{0.1cm}
  \begin{columns}[T]
    \begin{column}{0.31\textwidth}
      \begin{tcolorbox}[colback=white,colframe=pucpDarkGray,title=\textbf{CMD (Símbolo de Sistema)},fonttitle=\bfseries\tiny,boxrule=0.5pt,top=1.2mm,bottom=1.2mm]
        \tiny
        \begin{itemize}
          \item Intérprete clásico de MS-DOS.
          \item Sintaxis limitada para scripts modernos.
        \end{itemize}
      \end{tcolorbox}
    \end{column}

    \begin{column}{0.33\textwidth}
      \begin{tcolorbox}[colback=pucpLightBlue!40,colframe=pucpBlue,title=\textbf{PowerShell (Recomendado)},fonttitle=\bfseries\tiny,boxrule=0.5pt,top=1.2mm,bottom=1.2mm]
        \tiny
        \begin{itemize}
          \item Shell moderno nativo en Windows 10/11.
          \item Gran soporte de variables y automatizaciones.
        \end{itemize}
      \end{tcolorbox}
    \end{column}

    \begin{column}{0.34\textwidth}
      \begin{tcolorbox}[colback=white,colframe=pucpMidBlue,title=\textbf{Git Bash / Linux Shell},fonttitle=\bfseries\tiny,boxrule=0.5pt,top=1.2mm,bottom=1.2mm]
        \tiny
        \begin{itemize}
          \item Estándar en UNIX/Linux y macOS.
          \item Clave en servidores y clusters.
        \end{itemize}
      \end{tcolorbox}
    \end{column}
  \end{columns}

  \vspace{0.15cm}
  \begin{alertblock}{¿Qué es la variable PATH del sistema?}
    \small
    Es la lista de rutas donde Windows busca ejecutables (\texttt{git}, \texttt{agy} o \texttt{pdflatex}). Si un comando no se reconoce, su ruta debe agregarse al PATH.
  \end{alertblock}
\end{frame}

% ==============================================================================
% SECCIÓN 4: EXTENSIÓN OVERLEAF WORKSHOP Y CONEXIÓN REMOTA
% ==============================================================================
\section[Overleaf en IDE]{Integración de Overleaf en VS Code}

\begin{frame}{¿Por qué conectar Overleaf con VS Code?}
  \begin{columns}[T]
    \begin{column}{0.48\textwidth}
      \begin{alertblock}{Restricciones del Plan Web Gratuito}
        \small
        \begin{itemize}
          \item Tiempos límite de compilación estrictos (\textit{timeout errors}).
          \item Bloqueos en documentos pesados con muchas figuras.
          \item Límite de colaboradores simultáneos.
          \item Dependencia absoluta de la conexión al navegador.
        \end{itemize}
      \end{alertblock}
    \end{column}

    \begin{column}{0.48\textwidth}
      \begin{exampleblock}{Ventajas de Overleaf Workshop}
        \small
        \begin{itemize}
          \item \textbf{Compilaciones ilimitadas} y mayor velocidad.
          \item Sincronización bidireccional inmediata con la nube.
          \item Capacidad de compilar localmente con motores propios.
          \item Edición fluida con atajos avanzados y visor integrado.
        \end{itemize}
      \end{exampleblock}
    \end{column}
  \end{columns}
\end{frame}

\begin{frame}{Protocolo de Conexión: Autenticación por Cookies}
  \begin{columns}[T]
    \begin{column}{0.48\textwidth}
      \begin{block}{Paso 1: Obtener la Cookie}
        \begin{enumerate}\scriptsize
          \item Abre tu proyecto en el navegador (Overleaf).
          \item Presiona \texttt{F12} (o Clic derecho $\to$ \textit{Inspeccionar}).
          \item Ve a la pestaña \textbf{Network (Red)} y recarga (\texttt{F5}).
          \item En el filtro de búsqueda, escribe \texttt{cookie}.
          \item Copia el valor de \texttt{overleaf\_session}.
        \end{enumerate}
      \end{block}
    \end{column}

    \begin{column}{0.48\textwidth}
      \begin{block}{Paso 2: Vincular en VS Code}
        \begin{enumerate}\scriptsize
          \item Instala la extensión \textbf{Overleaf Workshop}.
          \item En el panel lateral, clic en \textit{Add New Server}.
          \item Selecciona \textit{Login with cookies}.
          \item Pega el valor de la cookie y presiona \texttt{Enter}.
          \item ¡Tus proyectos remotos aparecerán listados al instante!
        \end{enumerate}
      \end{block}
    \end{column}
  \end{columns}

  \vspace{0.15cm}
  \begin{tcolorbox}[colback=accentGreen!10,colframe=accentGreen,boxrule=0.5pt,arc=1.5mm,halign=center,top=1.5mm,bottom=1.5mm]
    \footnotesize\textbf{Resultado:} Acceso directo a \texttt{main.tex}, \texttt{referencias.bib} y figuras desde tu editor local.
  \end{tcolorbox}
\end{frame}

\begin{frame}{Flujo de Trabajo y Workspaces en VS Code}
  \begin{columns}[T]
    \begin{column}{0.50\textwidth}
      \begin{block}{Operaciones Clave en el Proyecto}
        \small
        \begin{itemize}
          \item \textbf{Compilar:} Botón \textit{Compile project} o atajo de teclado.
          \item \textbf{Visor PDF:} Panel lateral con \textit{View compiled PDF}.
          \item \textbf{Sincronización Automática:} Cambios guardados en VS Code se actualizan en Overleaf web.
          \item \textbf{Descarga Completa:} Clic derecho $\to$ \textit{Download} para respaldar todo localmente.
        \end{itemize}
      \end{block}
    \end{column}

    \begin{column}{0.48\textwidth}
      \begin{tcolorbox}[colback=pucpGray,colframe=pucpMidBlue,title=\textbf{Workspaces en VS Code},fonttitle=\bfseries\scriptsize,boxrule=0.6pt,top=1.5mm,bottom=1.5mm]
        \scriptsize
        \textbf{¿Por qué guardar un Workspace?}
        \begin{itemize}
          \item Genera un archivo \texttt{.code-workspace} con la sesión del proyecto.
          \item Permite reanudar el trabajo con un solo clic sin tener que reingresar credenciales.
          \item Menú: \textit{File} $\to$ \textit{Save Workspace As...}
        \end{itemize}
      \end{tcolorbox}
    \end{column}
  \end{columns}
\end{frame}

% ==============================================================================
% SECCIÓN 5: COPILOTOS DE IA Y GOOGLE ANTIGRAVITY
% ==============================================================================
\section[Copilotos \& IA]{Copilotos de IA y Google Antigravity}

\begin{frame}{Ecosistema de Modelos e Inteligencia Artificial}
  \begin{block}{Distinción Fundamental: Modelos vs. Asistentes/Agentes}
    \small
    Los \textbf{Modelos Fundacionales} (Gemini, Claude, GPT) generan texto; los \textbf{Agentes de Programación} (Antigravity, Claude Code, Cursor) pueden interactuar con tus archivos, terminales y proyectos.
  \end{block}

  \vspace{0.1cm}
  \begin{columns}[T]
    \begin{column}{0.31\textwidth}
      \begin{tcolorbox}[colback=white,colframe=pucpDarkGray,title=\textbf{OpenAI GPT / Codex},fonttitle=\bfseries\tiny,boxrule=0.5pt,top=1.2mm,bottom=1.2mm]
        \tiny
        \begin{itemize}
          \item Pionero en generación de código.
          \item Límites de cuota estrictos en versiones base.
          \item Menor integración en flujos locales libres.
        \end{itemize}
      \end{tcolorbox}
    \end{column}

    \begin{column}{0.31\textwidth}
      \begin{tcolorbox}[colback=white,colframe=pucpDarkGray,title=\textbf{Anthropic Claude Code},fonttitle=\bfseries\tiny,boxrule=0.5pt,top=1.2mm,bottom=1.2mm]
        \tiny
        \begin{itemize}
          \item Gran capacidad de razonamiento lógico.
          \item Mayor consumo de costos API.
          \item Consumo rápido de cuota de tokens.
        \end{itemize}
      \end{tcolorbox}
    \end{column}

    \begin{column}{0.34\textwidth}
      \begin{tcolorbox}[colback=pucpLightBlue!40,colframe=pucpRed,title=\textbf{Google Antigravity},fonttitle=\bfseries\tiny,boxrule=0.5pt,top=1.2mm,bottom=1.2mm]
        \tiny
        \begin{itemize}
          \item \textbf{Modo Agente:} Lee y edita archivos con confirmación.
          \item \textbf{Eficiente:} Gran ventana de contexto y bajo coste.
          \item Modelos flexibles: Gemini Flash/Pro y Claude.
        \end{itemize}
      \end{tcolorbox}
    \end{column}
  \end{columns}
\end{frame}

\begin{frame}{Capacidades de Google Antigravity en la Práctica}
  \begin{columns}[T]
    \begin{column}{0.48\textwidth}
      \begin{block}{Modos de Interacción}
        \small
        \begin{itemize}
          \item \textbf{Agent Mode:} Ejecuta tareas de forma autónoma (crear archivos, depurar \LaTeX, compilar).
          \item \textbf{Ask Mode:} Resuelve consultas conceptuales sin realizar cambios en tus archivos.
          \item \textbf{Plan Mode:} Diseña un plan de trabajo estructurado antes de aplicar modificaciones.
        \end{itemize}
      \end{block}
    \end{column}

    \begin{column}{0.48\textwidth}
      \begin{block}{Resolución Autónoma de Problemas}
        \small
        \begin{itemize}
          \item \textbf{Diagnóstico de Entorno:} Detecta y repara errores de comandos o rutas en el \texttt{PATH}.
          \item \textbf{Conversión de Formatos:} Transforma borradores Word/PDF a sintaxis limpia de \LaTeX.
          \item \textbf{Depuración Bibliográfica:} Normaliza entradas BibTeX y repara citas rotas.
        \end{itemize}
      \end{block}
    \end{column}
  \end{columns}
\end{frame}

\begin{frame}{Automatización de Git/GitHub con Antigravity}
  \begin{alertblock}{Control de Versiones mediante Lenguaje Natural}
    \small
    No necesitas memorizar la sintaxis de consola (\texttt{git add}, \texttt{commit}, \texttt{push}). Antigravity traduce tu petición en acciones atómicas de Git.
  \end{alertblock}

  \vspace{0.1cm}
  \begin{columns}[T]
    \begin{column}{0.49\textwidth}
      \textbf{\small Ejemplo de Prompt en Antigravity:}
      \begin{tcolorbox}[colback=pucpGray,colframe=pucpBlue,arc=1.5mm,boxrule=0.5pt,top=1.5mm,bottom=1.5mm]
        \scriptsize\ttfamily
        "Hola Antigravity, toma mi carpeta local de Overleaf y súbela a mi repositorio de GitHub dentro de 'prueba-vscode', haciendo commit y push a la rama main."
      \end{tcolorbox}
    \end{column}

    \begin{column}{0.49\textwidth}
      \textbf{\small Flujo Automatizado que ejecuta la IA:}
      \begin{enumerate}\scriptsize
        \item Identifica archivos (\texttt{main.tex}, \texttt{.bib}, figuras).
        \item Prepara y registra los cambios en Git.
        \item Formula un mensaje descriptivo de commit.
        \item Solicita tu autorización de seguridad.
        \item Sube los cambios al repositorio remoto.
      \end{enumerate}
    \end{column}
  \end{columns}
\end{frame}

% ==============================================================================
% SECCIÓN 6: MODALIDADES DE EVALUACIÓN Y TAREAS
% ==============================================================================
\section[Evaluación]{Modalidades de Evaluación y Portafolio}

\begin{frame}{Opciones para la Evaluación Práctica del Taller}
  {\small Se ofrecen dos modalidades para la entrega del trabajo calificado:}

  \vspace{0.1cm}
  \begin{columns}[T]
    \begin{column}{0.48\textwidth}
      \begin{block}{Opción 1: Modalidad Tradicional}
        \small
        \begin{itemize}
          \item Redacción de un \textbf{manuscrito académico} en \LaTeX.
          \item \textbf{Y además}, elaboración de una presentación en \textbf{Beamer} con las diapositivas.
          \item Entrega convencional mediante enlace de Overleaf o archivo ZIP.
        \end{itemize}
      \end{block}
    \end{column}

    \begin{column}{0.48\textwidth}
      \begin{exampleblock}{Opción 2: Modalidad Avanzada (Recomendada)}
        \small
        \begin{itemize}
          \item Desarrollar \textbf{solo uno de los dos productos}: Manuscrito \textbf{O} Presentación Beamer.
          \item Subir el trabajo mediante una rama (\textbf{Branch}) al repositorio colaborativo en GitHub.
          \item Automatizar la subida utilizando \textbf{Antigravity}.
        \end{itemize}
      \end{exampleblock}
    \end{column}
  \end{columns}

  \vspace{0.15cm}
  \begin{center}
    \footnotesize\textbf{Fecha límite de entrega:} Domingo al finalizar la semana del taller.
  \end{center}
\end{frame}

\begin{frame}{Beneficios Profesionales de la Modalidad Avanzada}
  \begin{center}
    \textbf{\large "Matar 5 pájaros de un solo tiro"}
  \end{center}
  \vspace{0.05cm}

  \begin{columns}[T]
    \begin{column}{0.31\textwidth}
      \begin{tcolorbox}[colback=pucpLightBlue!50,colframe=pucpBlue,title=\textbf{1. Menor Carga},fonttitle=\bfseries\scriptsize,halign=center,boxrule=0.5pt,top=1.5mm,bottom=1.5mm]
        \scriptsize
        Desarrollas un solo producto de alta calidad (Manuscrito o Beamer) en vez de duplicar esfuerzo.
      \end{tcolorbox}
    \end{column}

    \begin{column}{0.31\textwidth}
      \begin{tcolorbox}[colback=pucpLightBlue!50,colframe=pucpBlue,title=\textbf{2. Dominio de Git},fonttitle=\bfseries\scriptsize,halign=center,boxrule=0.5pt,top=1.5mm,bottom=1.5mm]
        \scriptsize
        Adquieres experiencia real en control de versiones y colaboración en GitHub.
      \end{tcolorbox}
    \end{column}

    \begin{column}{0.31\textwidth}
      \begin{tcolorbox}[colback=pucpLightBlue!50,colframe=pucpRed,title=\textbf{3. Portafolio LinkedIn},fonttitle=\bfseries\scriptsize,halign=center,boxrule=0.5pt,top=1.5mm,bottom=1.5mm]
        \scriptsize
        Evidencia verificable y pública de habilidades en \LaTeX, Overleaf, GitHub e IA.
      \end{tcolorbox}
    \end{column}
  \end{columns}

  \vspace{0.15cm}
  \begin{alertblock}{Construcción de Identidad Académica y Profesional}
    \small
    Un perfil en GitHub con código ordenado demuestra competencia técnica y distinción para becas, posgrados y convocatorias laborales.
  \end{alertblock}
\end{frame}

% ==============================================================================
% SECCIÓN 7: SÍNTESIS Y PRÓXIMOS PASOS
% ==============================================================================
\section[Cierre]{Síntesis y Próximos Pasos}

\begin{frame}{Resumen de la Sesión y Próxima Clase}
  \begin{columns}[T]
    \begin{column}{0.48\textwidth}
      \begin{block}{Lo que aprendimos hoy (Sesión 3)}
        \small
        \begin{itemize}
          \item Fundamentos de control de versiones con Git y GitHub.
          \item VS Code como entorno centralizado de trabajo.
          \item Conexión remota con Overleaf Workshop.
          \item Asistencia y automatización con Google Antigravity.
        \end{itemize}
      \end{block}
    \end{column}

    \begin{column}{0.48\textwidth}
      \begin{exampleblock}{Lo que veremos mañana (Sesión 4)}
        \small
        \begin{itemize}
          \item \textbf{Diseño de Presentaciones con Beamer}:
          \begin{itemize}\footnotesize
            \item Estructura de diapositivas y bloques temáticos.
            \item Paletas cromáticas institucionales y temas visuales.
            \item Inserción de fórmulas, figuras y tablas dinámicas.
          \end{itemize}
        \end{itemize}
      \end{exampleblock}
    \end{column}
  \end{columns}

  \vspace{0.15cm}
  \begin{center}
    \textbf{\small ¡Muchas gracias por su atención y participación!}\\[0.2ex]
    \scriptsize Dudas y asesorías disponibles por correo y en las sesiones complementarias de consulta.
  \end{center}
\end{frame}

\end{document}
